\section{Kết quả thực nghiệm}\label{sec:results}

Tất cả số liệu là trên \textbf{tập test}, sau khi chọn epoch (early stopping trên \code{val\_us}) và trọng số lớp $w$ (trên
validation). Với cấu hình chạy 3 seed (42, 43, 44), bảng ghi trung bình $\pm$ độ lệch chuẩn của macro F1; các cột khác lấy
trung bình. Cột ``Ep.'' là \emph{epoch tốt nhất / số epoch đã chạy} (early stopping dừng sau 10 epoch không cải thiện).

\subsection{Biểu diễn A --- mỗi đặc trưng là một bước (chia ngẫu nhiên, so sánh với Bài 1)}

\input{\tabdir/tab_A_rnn}
\input{\tabdir/tab_A_cnn}

\begin{figure}[H]\centering
\includegraphics[width=0.82\linewidth]{bars_A.png}
\caption{Biểu diễn A: macro F1 trên test của mọi cấu hình RNN/CNN/MLP và các mô hình của Bài 1 (cùng cách chia, cùng tiền xử lý).
Vạch đen: độ lệch chuẩn qua 3 seed.}\label{fig:bars_A}
\end{figure}

\paragraph{Nhận xét.}
\begin{itemize}
\item Mọi cấu hình RNN nằm trong khoảng macro F1 0,51--0,54, CNN trong khoảng 0,50--0,53. Tốt nhất là \code{A\_bigru}, \code{A\_gru\_2l}, \code{A\_lstm\_2l}
  ($\approx$0,535--0,537) và \code{A\_cnn\_bn} (0,528). Độ lệch chuẩn giữa các seed là 0,002--0,007, nên phần lớn
  khác biệt giữa các biến thể RNN (LSTM/GRU, 1/2 tầng, một/hai chiều, last/mean/attention) \textbf{không có ý nghĩa thống kê}.
\item \textbf{MLP} (không coi dữ liệu là chuỗi) đạt $\approx$0,52 --- ngang RNN/CNN. Việc ``xếp 11 đặc trưng thành chuỗi''
  không mang lại thông tin mới: thứ tự các đặc trưng là tuỳ ý, phép tích chập trên ``đặc trưng kề nhau'' và trạng thái ẩn ``truyền
  từ đặc trưng này sang đặc trưng kia'' không có ý nghĩa vật lý. Đúng như đề bài cảnh báo ở mục 14.
\item \textbf{XGBoost của Bài 1 (0,575) vẫn tốt hơn mọi mạng nơ-ron trên biểu diễn A} (kể cả Random Forest 0,555 cũng hơn).
  Điều này khớp với kết quả đã biết: với dữ liệu bảng có đặc trưng lệch phân phối mạnh và quan hệ dạng ngưỡng, mô hình cây vẫn mạnh
  hơn mạng nơ-ron [9]. Mạng nơ-ron chỉ vượt Logistic Regression (0,347) và KNN (0,416).
\item Learning rate nhỏ ($3\cdot10^{-4}$) cho kết quả kém nhất ở cả RNN và CNN (underfit, xem mục~\ref{sec:analysis});
  batch 2048 làm số bước cập nhật mỗi epoch giảm 4 lần, mô hình GRU chạm giới hạn 100 epoch trước khi hội tụ.
\item Với CNN, thêm BatchNorm giúp hội tụ nhanh và cho kết quả tốt nhất; Dropout (0,3) và GAP làm giảm kết quả --- mô hình
  $\approx$23 nghìn tham số trên 126 nghìn mẫu không thừa năng lực, điều chuẩn thêm chỉ làm nó underfit. Flatten tốt hơn GAP vì
  ở biểu diễn A \emph{vị trí} chính là \emph{danh tính của đặc trưng}.
\end{itemize}

\subsection{Biểu diễn B --- lịch sử của địa chỉ (chia theo địa chỉ)}

\input{\tabdir/tab_B}
\input{\tabdir/tab_group_other}

\begin{figure}[H]\centering
\includegraphics[width=0.78\linewidth]{bars_B.png}
\caption{Cùng cách chia theo địa chỉ: biểu diễn B, các biến thể đối chứng, XGBoost với 3 mức thông tin lịch sử và biểu diễn A.}
\label{fig:bars_B}
\end{figure}

\paragraph{Nhận xét.}
\begin{itemize}
\item \textbf{Biểu diễn B tăng macro F1 từ $\approx$0,53 lên $\approx$0,72--0,74} so với chính các mô hình đó trên biểu diễn A
  (cùng cách chia theo địa chỉ, Bảng~\ref{tab:group_other}). Đây là thông tin thật sự mới mà một dòng đơn lẻ không có: một địa
  chỉ ransomware (CryptoLocker, CryptoWall) nhận tiền nhiều ngày liên tiếp với mẫu hành vi lặp lại.
\item \textbf{Mô hình nơ-ron tốt nhất là CNN}: \code{B\_cnn\_k5} (0,742) và \code{B\_cnn\_3l} (0,737); RNN tốt nhất là
  \code{B\_gru\_2l} (0,735 $\pm$ 0,004). Trên B, \emph{tăng trường nhìn} (kernel 5, thêm tầng) có ích vì chuỗi dài hơn
  (tới 16 bước) và mẫu hành vi trải dài nhiều lần xuất hiện.
\item \textbf{Phần lợi ích do chuỗi mang lại}: \code{B\_gru\_L1} (chỉ bước hiện tại, nhưng có cờ ``lần đầu'' và khoảng cách ngày)
  đạt 0,660; L=4 đạt 0,712 và L=16 đạt 0,720. So với biểu diễn A trên cùng cách chia ($\approx$0,53), riêng việc \emph{biết địa chỉ
  từng xuất hiện} đã cho +0,13, và việc \emph{nhìn thấy nội dung các lần xuất hiện trước} cho thêm +0,06--0,08. L=32 không tốt hơn (0,692; dừng sớm ở epoch 26):
  địa chỉ có $>$16 lần xuất hiện rất ít, phần đệm dài chỉ thêm nhiễu.
\item \textbf{Đối chứng công bằng với XGBoost}: XGBoost chỉ dùng dòng hiện tại đạt 0,565; thêm 3 đặc trưng đếm
  (số lần đã xuất hiện, cờ lần đầu, khoảng cách ngày) đạt 0,759; thêm trung bình/max lịch sử đạt \textbf{0,767 --- tốt nhất toàn
  bộ đồ án}. Khi được cho cùng thông tin lịch sử dưới dạng đặc trưng tổng hợp thủ công, mô hình cây vẫn hơn RNN/CNN $\approx$0,025.
\item Biến thể \code{\_cnt} (thêm chính 3 đặc trưng đếm vào đầu vào phụ) không giúp GRU (0,722 so với 0,720, trong độ dao động seed) và
  không giúp CNN (0,731 so với 0,731): RNN/CNN \emph{đã tự học được} thông tin đếm từ chuỗi. Ngược lại, MLP cần nó (0,655 $\to$ 0,712).
\item LSTM và GRU cho kết quả gần nhau, chênh lệch $\le$0,01 và không nhất quán về chiều (A 1 tầng: LSTM hơn 0,009; A 2 tầng:
  bằng nhau; B: GRU hơn 0,005) --- không có loại nào vượt trội rõ ràng. Bi-GRU không giúp ở B: bước cuối
  (ngày hiện tại) là quan trọng nhất và đã được RNN một chiều đọc sau cùng. Mean pooling kém hơn hidden cuối vì làm loãng bước
  hiện tại.
\end{itemize}

\subsection{Kết quả theo độ dài lịch sử}

Bảng~\ref{tab:by_len} tách tập test theo số lần địa chỉ đã xuất hiện (tính cả ngày hiện tại).

\input{\tabdir/tab_by_len}

\begin{itemize}
\item Với địa chỉ \textbf{đã xuất hiện $\ge$ 3 lần}, RNN/CNN có lịch sử vượt XGBoost-lịch sử: ví dụ nhóm 6--15 lần
  F1 0,943 (\code{B\_cnn\_k5}) so với 0,885; nhóm $>$15 lần 0,973 (\code{B\_gru}) so với 0,937. Mô hình chuỗi khai thác được
  \emph{chi tiết từng lần xuất hiện} mà mean/max thủ công làm mất.
\item Với địa chỉ \textbf{xuất hiện lần đầu} (90\% số dòng test, 53\% số dòng ransomware), không có lịch sử nào để học, và
  XGBoost tốt hơn rõ (0,648 so với $\approx$0,55): mô hình cây khai thác đặc trưng của dòng hiện tại tốt hơn mạng nơ-ron
  (đúng như kết quả ở biểu diễn A). Vì nhóm này chiếm đa số, XGBoost thắng trên toàn bộ tập test.
\item Đường \code{B\_gru\_L1}/\code{B\_mlp} (không thấy nội dung lịch sử) tăng theo số lần xuất hiện chỉ nhờ cờ ``lần đầu'' và
  $\Delta$ngày; \code{B\_mlp\_cnt} (biết số lần đã xuất hiện) cao hơn, và RNN/CNN thấy toàn bộ chuỗi cao nhất ở các nhóm dài.
\end{itemize}

\subsection{F1 từng lớp và ma trận nhầm lẫn}

\input{\tabdir/tab_per_class}

\begin{figure}[H]\centering
\begin{subfigure}{0.32\linewidth}\includegraphics[width=\linewidth]{cm_A_random_A_bigru.png}\caption{RNN tốt nhất (A)}\end{subfigure}
\begin{subfigure}{0.32\linewidth}\includegraphics[width=\linewidth]{cm_B_group_B_cnn_k5.png}\caption{CNN tốt nhất (B)}\end{subfigure}
\begin{subfigure}{0.32\linewidth}\includegraphics[width=\linewidth]{cm_baseline_group_XGB_hist.png}\caption{XGBoost + lịch sử}\end{subfigure}
\caption{Ma trận nhầm lẫn trên tập test (chuẩn hoá theo hàng = recall từng lớp). Các ma trận khác nằm trong thư mục
\code{outputs/final/figures/} và notebook.}\label{fig:cm}
\end{figure}

\begin{itemize}
\item Lịch sử địa chỉ cải thiện mạnh các họ \textbf{tái sử dụng địa chỉ}: CryptoLocker 0,26 $\to$ 0,70, CryptoWall 0,38 $\to$ 0,84,
  otherRansom 0,08 $\to$ 0,66 (CNN-B; ở lớp này CNN-B còn vượt XGBoost-lịch sử 0,536).
\item Các họ \textbf{dùng địa chỉ mới mỗi lần} (Cerber, Locky: 99\% dòng là lần đầu xuất hiện) không được lợi gì từ lịch sử;
  ở đây XGBoost hơn mạng nơ-ron (Cerber 0,63 so với 0,52--0,56; Locky 0,745 so với 0,65--0,67).
\item Ma trận nhầm lẫn cho thấy lỗi \textbf{gần như hoàn toàn là nhầm giữa ransomware và white}; nhầm giữa các họ ransomware với
  nhau rất hiếm (55 trên 7.911 dòng ransomware với \code{B\_cnn\_k5}). Khó khăn của bài toán là \emph{phát hiện}, không phải
  \emph{phân biệt họ}.
\end{itemize}

\paragraph{Phân tích dự đoán sai (\code{B\_cnn\_k5}, chia theo địa chỉ).} Mô hình bỏ sót 2.694/7.911 dòng ransomware (34\%), trong
đó 78\% là địa chỉ xuất hiện lần đầu; 82\% trong 1.634 báo động nhầm cũng rơi vào địa chỉ lần đầu. Tách theo họ: với CryptoLocker,
CryptoWall và otherRansom, recall trên các dòng \emph{lặp lại} là 0,69/0,92/0,72 nhưng trên các dòng \emph{lần đầu} gần bằng 0 ---
mạng nơ-ron gần như chỉ nhận ra các họ này qua lịch sử, trong khi XGBoost-lịch sử vẫn bắt được 26\%/11\%/11\% dòng lần đầu nhờ đặc
trưng của chính giao dịch. Notebook (mục 10) liệt kê các dòng sai với độ tin cậy cao nhất và ví dụ chuỗi lịch sử tương ứng.

\subsection{Tốc độ huấn luyện}

Thời gian/epoch trong các bảng trên đo khi 6--8 tiến trình huấn luyện chạy song song trên cùng GPU, nên chỉ mang tính tham khảo.
Bảng~\ref{tab:bench} đo lại \emph{riêng từng mô hình} trên GPU rảnh.

\IfFileExists{\tabdir/tab_bench.tex}{\input{\tabdir/tab_bench}}{\textit{(Bảng đo tốc độ sẽ được sinh bởi \code{final/benchmark.py}.)}}

\begin{itemize}
\item \textbf{MLP nhanh nhất} (0,29--0,33 s/epoch): không có chiều chuỗi.
\item Ở biểu diễn A (chuỗi chỉ 11 bước, độ dài cố định), RNN 1 tầng và CNN 2 lớp có tốc độ \emph{gần như nhau}
  (0,39--0,43 s/epoch): cuDNN tối ưu RNN rất tốt cho chuỗi ngắn, còn chi phí BatchNorm/Dropout làm CNN chậm thêm (0,55 s).
  Thêm tầng làm chậm cả hai (LSTM 2 tầng 0,65 s, CNN 3 lớp 0,73 s).
\item Ở biểu diễn B (16 bước, độ dài thay đổi), \textbf{CNN nhanh hơn RNN} $\approx$25\% (0,68 so với 0,89--0,92 s/epoch):
  tích chập tính song song trên mọi vị trí, còn RNN phải tính tuần tự và xử lý \code{pack\_padded\_sequence}. Bi-GRU chậm nhất
  (2,14 s/epoch, gấp 2,4 lần GRU một chiều) mà không tốt hơn.
\item \textbf{GRU ít tham số hơn LSTM} (3 cổng so với 4: 22.439 so với 27.687 tham số ở cấu hình gốc A) và nhanh hơn một chút
  (0,39 so với 0,41 s; 0,53 so với 0,65 s khi 2 tầng).
\item Như vậy ở biểu diễn B, CNN vừa \textbf{nhanh hơn} vừa \textbf{tốt hơn} RNN.
\end{itemize}

\subsection{Trả lời các câu hỏi của đề}

\begin{itemize}
\item \textbf{Mô hình tốt nhất?} Trong các mạng nơ-ron: \code{B\_cnn\_k5} (CNN 1 chiều trên lịch sử địa chỉ, macro F1 0,742).
  Trên toàn bộ đồ án: XGBoost với đặc trưng lịch sử tổng hợp (0,767). Với cùng thông tin một dòng (biểu diễn A), XGBoost Bài 1
  (0,575) tốt hơn mọi RNN/CNN ($\le$0,54).
\item \textbf{Vì sao?} Yếu tố quyết định là \emph{thông tin đầu vào}, không phải loại mạng: thêm lịch sử địa chỉ tăng
  $\approx$0,2 macro F1, còn mọi thay đổi kiến trúc trong cùng một biểu diễn chỉ thay đổi $\le$0,03. CNN tốt hơn RNN một chút trên B vì
  mẫu ``nhận tiền lặp lại vài ngày liên tiếp'' là mẫu \emph{cục bộ}, đúng điểm mạnh của tích chập; XGBoost tốt nhất vì phần lớn
  dòng test không có lịch sử và ở đó mô hình cây dùng đặc trưng dạng bảng hiệu quả hơn.
\item \textbf{Mô hình nhanh nhất?} MLP (Bảng~\ref{tab:bench}). Trong các mô hình chuỗi: ở A, GRU/LSTM 1 tầng và CNN đơn giản
  ngang nhau ($\approx$0,4 s/epoch); ở B, CNN nhanh hơn RNN $\approx$25\%.
\item \textbf{Mô hình ổn định nhất?} Độ lệch chuẩn qua 3 seed của các cấu hình đều nhỏ (0,001--0,008; riêng \code{B\_gru\_cnt} 0,015), nhỏ hơn nhiều so với
  khoảng cách A--B. CNN có BN + Dropout trên B cho đường validation mượt và khoảng cách train--val nhỏ nhất (mục~\ref{sec:analysis}).
\item \textbf{Kết quả có hợp lý không?} Có. Accuracy luôn $>$98\% nhưng vô nghĩa (đoán toàn \code{white} đã đạt 98,6\%), nên
  macro F1 được dùng làm độ đo chính. Lớp \code{white} luôn có F1 $\approx$0,99; các họ có nhiều địa chỉ lặp lại được cải thiện mạnh
  nhờ lịch sử, các họ dùng địa chỉ một lần thì không --- đúng với phân tích dữ liệu ở mục 2. Cách chia theo địa chỉ đảm bảo kết
  quả của B không đến từ việc ``nhớ'' địa chỉ đã thấy khi huấn luyện.
\end{itemize}
